\section{Теоретический результат: область устойчивости метода с инерцией}\label{sec:theory}

\subsection{Двумерное линейное отображение}

Рассмотрим скалярную квадратичную функцию потерь
\begin{equation}
  L(w) = \tfrac12 \lambda w^2, \qquad \lambda > 0,
\end{equation}
и градиентный спуск с инерцией в форме~\eqref{eq:hb}, где $\nabla L(w) = \lambda w$:
\begin{equation}\label{eq:hb2}
  v_{t+1} = \beta v_t + \lambda w_t, \qquad w_{t+1} = w_t - \eta v_{t+1}.
\end{equation}
Подставляя первое равенство во второе, получаем
$w_{t+1} = w_t - \eta(\lambda w_t + \beta v_t) = (1-\eta\lambda)w_t - \eta\beta v_t$.
Следовательно, пара $(w_t,v_t)$ подчиняется линейному отображению
\begin{equation}\label{eq:matrix}
  \begin{pmatrix} w_{t+1} \\ v_{t+1}\end{pmatrix}
  = M \begin{pmatrix} w_t \\ v_t\end{pmatrix},
  \qquad
  M = \begin{pmatrix} 1-\eta\lambda & -\eta\beta \\[2pt] \lambda & \beta \end{pmatrix}.
\end{equation}

\subsection{Условия устойчивости}

\begin{proposition}[Область устойчивости метода тяжёлого шарика]\label{prop:stab}
Пусть $x := \eta\lambda$. Нулевое решение отображения~\eqref{eq:matrix}
асимптотически устойчиво тогда и только тогда, когда
\begin{equation}\label{eq:region}
  |\beta| < 1 \qquad\text{и}\qquad 0 < x < 2(1+\beta).
\end{equation}
При $\beta = 0$ условие переходит в классическое условие устойчивости
градиентного спуска $0 < \eta\lambda < 2$.
\end{proposition}

\begin{proof}
Вычислим инварианты матрицы $M$:
\begin{equation}
  \operatorname{tr} M = 1 - \eta\lambda + \beta = 1 + \beta - x,
\end{equation}
\begin{equation}
  \det M = (1-\eta\lambda)\beta + \eta\beta\lambda = \beta - \eta\lambda\beta + \eta\lambda\beta = \beta .
\end{equation}
Существенно, что определитель не зависит от шага обучения и равен коэффициенту
инерции; это отражает тот факт, что инерция задаёт скорость сжатия фазового
объёма.

Характеристический полином имеет вид $p(\mu) = \mu^2 - T\mu + \Delta$, где
$T=\operatorname{tr}M$, $\Delta = \det M = \beta$. Асимптотическая устойчивость
равносильна тому, что оба корня лежат строго внутри единичной окружности.
По критерию Жюри для полинома второй степени это эквивалентно системе
\begin{equation}
  |\Delta| < 1, \qquad p(1) > 0, \qquad p(-1) > 0,
\end{equation}
то есть
\begin{equation}
  |\beta| < 1, \qquad 1 - T + \beta > 0, \qquad 1 + T + \beta > 0 .
\end{equation}
Подставляя $T = 1+\beta-x$, получаем из второго неравенства
$1 - 1 - \beta + x + \beta = x > 0$, а из третьего
$1 + 1 + \beta - x + \beta = 2 + 2\beta - x > 0$, то есть $x < 2(1+\beta)$.
Вместе с $|\beta|<1$ это в точности~\eqref{eq:region}.
\end{proof}

\begin{corollary}[Характер сходимости]
Корни комплексны при $T^2 < 4\Delta$, то есть при
$(1+\beta-x)^2 < 4\beta$. В этом случае
\begin{equation}
  |\mu_{1,2}| = \sqrt{\Delta} = \sqrt{\beta},
\end{equation}
и скорость сходимости не зависит от $\lambda$: все собственные направления
затухают с одинаковым множителем $\sqrt{\beta}$ за итерацию. Это объясняет
осцилляторный характер траекторий при больших $\beta$ и то, почему инерция
выравнивает скорости сходимости вдоль направлений с разной кривизной.
\end{corollary}

\begin{corollary}[Расширение допустимого шага]
Максимальный устойчивый шаг равен $\eta_{\max}(\beta) = 2(1+\beta)/\lambda$,
что при $\beta\to 1$ вдвое превышает соответствующее значение для градиентного
спуска. Порог грани устойчивости для метода с инерцией, таким образом, есть
$\lambda_{\max}(\nabla^2 L) \approx 2(1+\beta)/\eta$.
\end{corollary}

\subsection{Численная проверка}

На рис.~\ref{fig:theory} слева показан спектральный радиус $\rho(M)$ в
плоскости $(x,\beta)$ с нанесённой аналитической границей $x = 2(1+\beta)$;
справа --- множество пар, для которых численная итерация~\eqref{eq:hb2} из
начального условия $w_0=1$, $v_0=0$ за 400 шагов даёт $|w_{400}| < 10^{-3}$.
Численная область устойчивости совпадает с аналитической с точностью до
разрешения сетки, что подтверждает предложение~\ref{prop:stab}. Вертикальная
пунктирная линия $x=2$ отмечает порог обычного градиентного спуска: вся область
правее неё и левее красной кривой --- это шаги, устойчивые только благодаря
инерции.

\begin{figure}[htbp]
  \centering
  \includegraphics[width=\linewidth]{f05_theory_stability.png}
  \caption{Область устойчивости метода тяжёлого шарика на квадратичной функции.
  Слева --- спектральный радиус матрицы перехода (значения обрезаны сверху
  на уровне 2); белая линия --- уровень $\rho=1$, красная штриховая ---
  аналитическая граница $\eta\lambda = 2(1+\beta)$. Справа --- численно
  определённая область устойчивости.}
  \label{fig:theory}
\end{figure}

\subsection{Почему в нелинейном случае граница усложняется}

Предложение~\ref{prop:stab} даёт гладкую границу: в плоскости $(x,\beta)$ это
прямая. Однако три обстоятельства разрушают эту гладкость при переходе к
нейронной сети.

Во-первых, в многомерном случае $L$ имеет спектр кривизн
$\lambda_1,\dots,\lambda_d$, и условие устойчивости должно выполняться
одновременно для всех мод: $\eta\lambda_i < 2(1+\beta)$. Граница области
устойчивости определяется максимальным собственным значением гессиана, которое
само зависит от точки в пространстве параметров, а значит --- от всей
предыстории траектории.

Во-вторых, при обучении нейронной сети гессиан не постоянен: система, вышедшая
за порог устойчивости по одной моде, не обязана расходиться, поскольку может
перейти в область меньшей кривизны (catapult-механизм~\cite{lewkowycz2020}) и
стабилизироваться на грани устойчивости~\cite{cohen2021}. Итоговый исход
зависит от того, успеет ли траектория попасть в такую область раньше, чем
норма весов достигнет переполнения; это делает исход чувствительным к сколь
угодно малым изменениям гиперпараметров.

В-третьих, при потере устойчивости неподвижной точки нелинейная система не
уходит на бесконечность сразу, а проходит через каскад удвоения
периода~\cite{feigenbaum1978} и хаос. На рис.~\ref{fig:bif} показана
бифуркационная диаграмма для скалярной нелинейной задачи
$L(w) = \tfrac12 (w^2-y)^2$ при трёх значениях инерции. Видно, что рост
$\beta$ смещает каскад бифуркаций в область меньших $\eta$ и делает переход
более резким: при $\beta=0.9$ окно устойчивого поведения по $\eta$ существенно
сужается, хотя для квадратичной модели инерция допустимый шаг, наоборот,
расширяет. Это принципиальное различие между линейным и нелинейным случаем и
объясняет, почему граница, гладкая в теории для квадратики, становится
изрезанной для сети.

\begin{figure}[htbp]
  \centering
  \includegraphics[width=\linewidth]{f06_bifurcation.png}
  \caption{Бифуркационная диаграмма для нелинейной скалярной задачи
  $L(w)=\tfrac12(w^2-y)^2$ при $y=1$: показаны последние 120 из 600 итераций
  для каждого $\eta$. Слева направо: $\beta=0$, $\beta=0.5$, $\beta=0.9$.}
  \label{fig:bif}
\end{figure}
